\documentclass[a4paper,11pt]{article}
\usepackage[top=1.5cm, bottom=0.5cm, left=1.5cm, right=2cm]{geometry}

\usepackage{graphicx}
% \usepackage{url}
% \usepackage{palatino}
\usepackage{fontspec}
\setmainfont{TeX Gyre Pagella}

\usepackage{array}
\usepackage{ragged2e}        % necesario para \RaggedRight en \newcolumntype
\usepackage{tabularx}
\usepackage{multirow}
\usepackage{enumitem}

% \usepackage{color}
\usepackage{xcolor}
\usepackage{hyperref}
\usepackage{yaac-icons}
\usepackage[table]{xcolor}
\usepackage{amssymb}
\usepackage{booktabs}


\usepackage{multirow}

% \newcolumntype{L}[1]{>{\RaggedRight\arraybackslash}p{#1}}

\newcommand{\topbottomstrut}{\rule{0pt}{1.1em}} % “strut” para altura interna
\setlength{\extrarowheight}{0pt}
\renewcommand\arraystretch{1}

\usepackage[framemethod=tikz]{mdframed} % dibuja marcos finos y limpios

% \textheight=9.75in
\raggedbottom
\setlength{\tabcolsep}{0in}
\renewcommand{\labelitemii}{$\circ$}
\setlength{\parindent}{0pt}
\setlength{\tabcolsep}{0pt}
\renewcommand\arraystretch{1}

\usepackage{array,graphicx,multirow,ragged2e}
\usepackage[table]{xcolor} % lädt colortbl mit

\definecolor{background}{RGB}{216,228,236}
\definecolor{mygrey}{gray}{0.75}
\definecolor{myblue}{RGB}{0,40,70}
\definecolor{lightgrey}{RGB}{240,240,240}
\definecolor{accent}{RGB}{0,140,200} % azul acento para el degradado/bandas finas
\definecolor{qrblue}{RGB}{0,42,68} 
\definecolor{cvbackgroundgray}{RGB}{235,235,235}
\definecolor{linksoft}{RGB}{30,100,160} % azul más claro
\definecolor{linksoft2}{RGB}{100,160,220}

\colorlet{symbolcolor}{linksoft}
\colorlet{linkcolor}{linksoft}

\hypersetup{
  colorlinks = true,
  linkcolor  = linksoft,
  urlcolor   = linksoft,
  citecolor  = linksoft,
}

% \hypersetup{
%   colorlinks = true,
%   linkcolor  = linksoft,
%   urlcolor   = linksoft,
%   citecolor  = linksoft,
%   pdftitle={Manuel Caipo - CV},
%   pdfauthor={Manuel Caipo},
%   pdfsubject={Curriculum Vitae},
%   unicode=true,
%   bookmarksopen=false,
%   pdfstartview=FitH,
%   pdfpagemode=UseNone,
%   breaklinks=true
% }

\usepackage{fancyhdr}
\usepackage{lastpage} % solo si quieres “Página X de Y”

\pagestyle{empty}


\newmdenv[
  linecolor=mygrey,           % o black
  linewidth=0.6pt,
  leftline=true,
  rightline=false, topline=false, bottomline=false,
  innerleftmargin=8pt,
  innerrightmargin=0pt,
  innertopmargin=0pt,         % <-- sin acolchado arriba
  innerbottommargin=0pt,
  skipabove=0pt, skipbelow=0pt
]{expblock}



\newcolumntype{L}[1]{>{\RaggedRight\arraybackslash}p{#1}} % texto pegado izquierda
\newcolumntype{R}[1]{>{\RaggedLeft\arraybackslash}p{#1}} 



% fecha con ajuste vertical opcional: \datecell[<desplazamiento>]{<texto>}
\newcommand{\datecell}[2][0pt]{%
  \raisebox{#1}{%
    \parbox[t]{\datecol}{\raggedleft\textbf{#2}\strut}%
  }%
}

% En el preámbulo (para viñetas cuadradas más pequeñas)
%\renewcommand{\labelitemi}{\tiny$\blacksquare$}



\usepackage{tikz}
\usetikzlibrary{calc,shadows.blur,decorations.pathmorphing,fadings}
\usetikzlibrary{shadings}
\usepackage{eso-pic}      % para dibujar sobre la página
\usepackage{qrcode}       % para generar el código QR


\usepackage{needspace} % para reservar espacio y evitar saltos


\newcommand{\isep}{-2 pt}
\newcommand{\lsep}{-0.5cm}
\newcommand{\lsepe}{-0.4cm}
\newcommand{\psep}{-0.6cm}
\newcommand{\resitem}[1]{\item #1 \vspace{-2pt}}

\newcommand{\sq}[1]{\raisebox{0.25ex}{\rule{#1}{#1}}}

% Tamaño recomendado: 0.45ex–0.70ex (prueba y elige)
\setlist[itemize]{%
  label=\sq{0.65ex},   % tamaño de la viñeta (más pequeña que 0.75ex)
  labelsep=0.4em,     % espacio viñeta–texto (reduce si quieres aún más: 0.3em)
  leftmargin=0.5em,    % sangría del bloque (reduce “espacio antes” de la viñeta)
  labelindent=0pt,     % empuje extra de la viñeta (déjalo 0 salvo que quieras mover solo el símbolo)
  itemsep=1pt,         % espacio entre ítems
  parsep=0pt,          % sin espacio entre párrafos internos
  topsep=0.05pt,          % espacio antes y después de la lista
  partopsep=0pt        % sin extra por empezar párrafo
}



% --- en el preámbulo ---
\newlength{\framepad}
\setlength{\framepad}{8pt} % el mismo padding que quieres para ambas cajas

% helper para cajas a todo el ancho, perfectamente alineadas
\newcommand{\FullColorBox}[2]{%
  \begingroup
  \setlength{\fboxsep}{\framepad}%
  \noindent\colorbox{#1}{%
    \parbox{\dimexpr\textwidth\fboxsep\relax}{#2}%
  }%
  \endgroup
}

% --- Definición de resheading ---
% \newcommand{\resheading}[1]{%
%   \begingroup
%   \setlength{\fboxsep}{6pt}%
%   \noindent\colorbox{background}{%
%     \parbox{\dimexpr\linewidth-0.01\fboxsep\relax}{%
%       \footnotesize\textbf{\textcolor{myblue}{#1}}%
%     }%
%   }%
%   \endgroup
% }

% --- Definición de resheading para CV (sin saltos de página) ---
\newcommand{\resheading}[1]{%
  \phantomsection % Crea el ancla para el link sin saltar de página
  \addcontentsline{toc}{section}{#1} % Agrega el título al árbol del PDF
  \begingroup
  \setlength{\fboxsep}{6pt}%
  \noindent\colorbox{background}{%
    \parbox{\dimexpr\linewidth-0.01\fboxsep\relax}{%
      \footnotesize\textbf{\textcolor{myblue}{#1}}%
    }%
  }%
  \endgroup
}



\newlength{\myPad}
\setlength{\myPad}{0pt} % padding que quieras

% Caja a todo el ancho con relleno
\newcommand{\fullpadbox}[2]{% #1=color, #2=contenido
  \begingroup
  \setlength{\fboxsep}{\myPad}% controla el padding
  \noindent
  \colorbox{#1}{%
    \parbox{\dimexpr\linewidth+3\fboxsep\relax}{#2}%
  }%
  \endgroup
}


% --- Topbar limpio, angosto y con QR a juego ---
\newcommand{\FancyTopbar}{%
  \begin{tikzpicture}[remember picture,overlay]
    % Alturas: más angosto y pegado arriba
    \def\hOne{6mm}      % altura franja principal (antes 8mm)
    \def\hTwo{12mm}     % alcance piezas diagonales (antes 16mm)
    \def\overshoot{0.8pt} % evita micro-hueco en el borde superior

    % Color del QR (ajústalo si quieres más contraste)
    \colorlet{qrcolor}{myblue!15!white}

    % Franja sólida superior (con overshoot)
    \fill[myblue]
      ($(current page.north west)+(0,\overshoot)$) rectangle
      ($(current page.north east)+(0,-\hOne)$);

    % Pieza diagonal clara (efecto geométrico, sin líneas)
    \path[fill=background]
      ($(current page.north east)+(0,-\hOne)$) --
      ($(current page.north east)+(-6.2cm,-\hOne)$) --
      ($(current page.north east)+(-4.6cm,-\hTwo)$) --
      ($(current page.north east)+(0,-\hTwo)$) -- cycle;

    % Transición sutil (puedes eliminarla si quieres 100% plano)
    \path[shade, left color=accent, right color=myblue, opacity=0.85]
      ($(current page.north east)+(-6.2cm,-\hOne)$) --
      ($(current page.north east)+(-10.5cm,-\hOne)$) --
      ($(current page.north east)+(-9.0cm,-\hTwo)$) --
      ($(current page.north east)+(-4.6cm,-\hTwo)$) -- cycle;


  \end{tikzpicture}%
}

\newcommand{\FancyQRBottom}{%
  \begin{tikzpicture}[remember picture,overlay]
    % color del QR: en sintonía con tus enlaces (linksoft) o con myblue
    \colorlet{qrcolor}{linksoft} % o: \colorlet{qrcolor}{myblue!20!white}

    % Nodo en la esquina inferior derecha, con pequeño margen interior
    \node[
      anchor=south east,
      inner sep=1mm,
      fill=white,
      draw=qrcolor,
      line width=0.45pt,
      rounded corners=1mm
    ] at ($(current page.south east)+(-4mm,4mm)$) {%
      \begingroup
        \color{qrcolor}
        \qrcode[height=10mm]{https://manuelcaipocc.github.io/}%
      \endgroup
    };
  \end{tikzpicture}%
}

\makeatletter
\AddToHook{shipout/foreground}{%
  \begin{tikzpicture}[remember picture,overlay]
    % Ajusta el + (0,8mm) si lo quieres más alto/bajo
    \node[anchor=south, inner sep=0pt] 
      at ($(current page.south)+(0,8mm)$)
      {\footnotesize \textcolor{myblue}{\thepage}};
  \end{tikzpicture}%
}

% ===== Macro para el enlace de certificado =====
\newcommand{\cert}[1]{\raisebox{0.30ex}{\href{#1}{\tiny [Z]}}}

% ===== Macro para instituciones con enlace =====
\newcommand{\inst}[2]{\href{#1}{#2}}
% --- Aplica en todas las páginas (o usa shipout/firstpage si solo la primera) ---
% \AddToHook{shipout/foreground}{\FancyTopbar}
% \AddToHook{shipout/foreground}{\FancyQRBottom}

\AddToHook{shipout/background}{\FancyTopbar}
\AddToHook{shipout/background}{\FancyQRBottom}

\makeatother
% ----------------------------------------------------------------

\begin{document}

\setlength{\tabcolsep}{6pt}
\begin{tabular}{L{0.7\linewidth} R{0.275\linewidth}}
  
  \vspace{-8em}
  {\Large\bfseries Manuel Caipo}

  
  \vspace{0.5em}
  \textbf{Machine Learning Engineer \,|\, Mechanical Engineer}

  \vspace{0.5em}
  % ==== Personal Data ====
  \footnotesize
  15.02.97, Peru \,|\, Schubartstraße 62 \,|\, 89134 Blaustein, Germany\par
  \href{mailto:manuelcaipocc@outlook.com}{manuelcaipocc@outlook.com} \,|\, \smartphone{+49 176 32618095}\par
  \github{https://github.com/manuelcaipocc}\,|\, 
  \linkedin{https://www.linkedin.com/in/manuel-caipo-89845b151}\,|\, 
  \personalwebsite{https://manuelcaipocc.github.io/}\,|\, 
  \workpermit{https://manuelcaipocc.github.io/cv-caipo-latex/cv_links/ID/Aufenthaltstitel_caipo.pdf}


\setlength{\tabcolsep}{3pt}
\begin{tabular}{@{}l l l@{}}
German: C1 \,[\href{https://manuelcaipocc.github.io/cv-caipo-latex/cv_links/german/testdaf-digital-certificate-scale.pdf}{Certificate}] &
English: C1 &
Spanish: Native
\end{tabular}

&
  \includegraphics[width=0.5\linewidth]{images/manuelcaipo.jpg}
\end{tabular}

\vspace{-0.4em}


\begin{tabular}{p{\linewidth}}
  \footnotesize
  \justifying
    \cellcolor{cvbackgroundgray}
\noindent Mechanical Engineer and Machine Learning Specialist focused on physics-aware modeling and AI-based condition monitoring for industrial cyber-physical systems. Strong ability to interpret sensor behavior, analyze physical degradation mechanisms, and define physically consistent anomaly detection criteria before implementing scalable ML solutions. Experienced in high-frequency industrial time-series modeling, data pipeline automation, and end-to-end system deployment. Demonstrated industrial impact while pursuing mathematically grounded approaches to predictive maintenance and structured system modeling.
 My recent work increasingly emphasizes structured and topology-aware system representations that encode physical relationships and constraints, reflecting a growing research-oriented focus on interpretable and physically grounded AI models for complex engineering architectures.
\end{tabular}



%  data engineer 
% Machine Learning Engineer with a strong data engineering focus, experienced in designing scalable data architectures and automating high-frequency industrial time-series pipelines for AI-driven predictive maintenance systems. Skilled in building robust ETL workflows, integrating shop-floor IoT streams with cloud-based analytics platforms (Azure ML, Snowflake, Databricks, OPC UA, MQTT), and enabling reliable end-to-end ML deployment. Demonstrated measurable business impact through improvements in asset availability, production output, and operational efficiency. Mechanical engineering background supporting physics-aware system design.


 
% machine learning 
% Machine Learning Engineer and Mechanical Engineer specializing in AI-driven predictive maintenance and physics-aware condition monitoring for industrial systems. Experienced in developing, validating, and deploying ML/DL models for anomaly detection and RUL prediction, supported by scalable data engineering pipelines for high-frequency time-series data. Strong integration of IoT data streams across shop-floor and cloud analytics layers (Azure ML, Snowflake, Databricks, OPC UA, MQTT). Proven business impact through measurable improvements in asset availability and production performance.

% research
% Mechanical Engineer and Machine Learning Specialist focused on physics-aware modeling and AI-based condition monitoring for industrial cyber-physical systems. Strong ability to interpret sensor behavior, analyze physical degradation mechanisms, and define physically consistent anomaly detection criteria before implementing scalable ML solutions. Experienced in high-frequency industrial time-series modeling, data pipeline automation, and end-to-end system deployment. Demonstrated industrial impact while pursuing mathematically grounded approaches to predictive maintenance and structured system modeling.
% My recent work increasingly emphasizes structured and topology-aware system representations that encode physical relationships and constraints, reflecting a growing research-oriented focus on interpretable and physically grounded AI models for complex engineering architectures.

% \noindent Mechanical Engineer and Machine Learning researcher working at the intersection of artificial intelligence and engineering systems. My work focuses on physics-aware machine learning for complex cyber-physical systems, combining industrial domain knowledge with data-driven models to develop interpretable and reliable AI methods.
% In industrial settings, I have developed scalable machine learning pipelines for high-frequency sensor data, addressing anomaly detection, degradation modeling, and remaining useful life prediction for large-scale hydraulic systems. This experience has motivated my research interest in structured learning approaches that explicitly incorporate system topology, physical constraints, and engineering design knowledge into machine learning models.
% My long-term research goal is to develop machine learning methods that integrate physical structure, topology, and system architecture into learning algorithms, enabling robust AI systems capable of reasoning about complex engineered systems.

\vspace{0.2em}


% ======= EXPERIENCE (same width as \textwidth) =======

\resheading{PROFESSIONAL EXPERIENCE}

\vspace{-0.25em}
\noindent
% --- Widths coupled to \textwidth ---
\newlength{\datecol}
\setlength{\datecol}{2.7cm}   % width of the date column (adjust if you want)

\newlength{\colgap}
\setlength{\colgap}{6mm}      % visual separation between date and content


\begin{tabular}{@{}>{\raggedleft\arraybackslash}m{\datecol} @{\hspace{\colgap}} p{\dimexpr\textwidth-\datecol-\colgap\relax}@{}}


% ===== Mercedes-Benz =====
\datecell[-1.6ex]{\small 06.2026 – Today} &
\begin{minipage}[t]{\linewidth}\vspace{0pt}
\begin{expblock}
  \footnotesize
  \setlength{\parindent}{0pt}\setlength{\parskip}{0pt}
  \textbf{\href{https://group.mercedes-benz.com}{Mercedes-Benz AG, Sindelfingen, Germany}}\\[-0.1em]
  \textbf{RD/EDD eDrive Software Integration – PhD Candidate}\\[-1em]
  \begin{itemize}[nosep,leftmargin=1em]
    \item Doctoral research, in cooperation with TUM's Chair of Sustainable Mobile Drivetrains, on context-aware multimodal learning for fault analysis in software-intensive eDrive systems.
    \item Design of agentic LLM approaches that orchestrate specialized ML fault classifiers across heterogeneous data sources (signals, logs, software and test artifacts).
    \item Focus on explainable and traceable diagnostics, distilling validated expert knowledge into verifiable diagnostic decisions to improve engineering processes.
  \end{itemize}
\end{expblock}
\end{minipage}
\\[-0.2em]

% ===== Bosch =====
\datecell[-1.6ex]{\small 11.2024 – 03.2026} &
\begin{minipage}[t]{\linewidth}\vspace{0pt}
\begin{expblock}
  \footnotesize
  \setlength{\parindent}{0pt}\setlength{\parskip}{0pt}
  \textbf{\href{https://www.boschrexroth.com/en/dc/}{Bosch Rexroth AG, Ulm, Germany}}\\[-0.1em]
  \textbf{DC-SSD3 – Master's Thesis }\\[-1em]
    \begin{itemize}[nosep,leftmargin=1em]
    \item Master's thesis focusing on the development of an ML/AI-based condition monitoring module for \emph{CytroConnect}, combining physical models, hydraulic domain knowledge, and data compliant with the \emph{Fluid 4.0} standard.
    \end{itemize}
  \textbf{DC-SSD3 – Data Scientist Working Student}\\[-1em]
  \begin{itemize}[nosep,leftmargin=1em]
\item Built scalable ETL pipelines for high-frequency industrial time-series data (SQL/NoSQL) and integrated message brokers (Solace/MQTT).
\item Developed and deployed ML/DL-based predictive models for anomaly detection and RUL prediction in hydraulic systems.
\item Applied explainable AI methods to increase transparency, robustness, and acceptance of data-driven models.

  \end{itemize}
\end{expblock}
\end{minipage}
\\[-0.2em]


% ===== Freeport =====
\datecell[-1.6ex]{\small 04.2021 – 07.2023} &
\begin{minipage}[t]{\linewidth}\vspace{0pt}
\begin{expblock}
  \footnotesize
  \setlength{\parindent}{0pt}\setlength{\parskip}{0pt}
  \textbf{\href{https://www.fcx.com}{Freeport-McMoRan Cerro Verde, Phoenix, U.S.A.}}\\[-0.1em]
  \textbf{Data Analytics – Junior Data Scientist (1 → 2)}\\[-1em]
\begin{itemize}[nosep,leftmargin=1em]
  \item Developed and deployed ML/DL models (LSTM, CNN, XGBoost) in Azure ML for anomaly detection, wear and RUL prediction of critical assets (FLSmidth primary crushers, Weir Minerals cyclone pumps), including SHAP/PDP explainability, continuous retraining and optimized SQL Stored Procedures for accelerated training data preparation.
  \item RUL-based proactive maintenance planning: +1.5\% asset availability (≈\$10M/year savings). Awarded Innova 2022, Showroom Award and President’s Award.
\end{itemize}

  \textbf{Maintenance Reliability – Trainee → Junior Data Analyst}\\[-1em]
  \begin{itemize}[nosep,leftmargin=1em]
\item Time-series analysis (ARIMA), RCA and automated KPI reporting (SQL, Dagster, Power BI) to shift maintenance from fixed intervals to condition-based strategies in SAP.
\item Long-term performance and failure analysis of the concentrator plant: primary and secondary crushers, HPGR, ball mills, conveyor belts and cyclone pumps.
  \end{itemize}
\end{expblock}
\end{minipage}
\\[-0.2em]

% ===== IMCO =====
\datecell[-1.6ex]{\small 08.2020 – 03.2021} &
\begin{minipage}[t]{\linewidth}\vspace{0pt}
\begin{expblock}
  \footnotesize
  \setlength{\parindent}{0pt}\setlength{\parskip}{0pt}
  \textbf{\href{https://www.linkedin.com/company/imco-servicios-s-a-c/posts/?feedView=all}{IMCO Servicios S.A.C, Arequipa, Peru}}\\[-0.1em]
  \textbf{Engineering and Development – Junior Engineer}\\[-1em]
  \begin{itemize}[nosep,leftmargin=1em]
    \item Performed static and dynamic structural analyses with topology optimization in a BIM-oriented methodology, including architectural modeling (Tekla), FEM-based structural design (SAP2000, IdeaStatica, Inventor, Ansys Structural), and final plan review in AutoCAD 3D.

  \end{itemize}
\end{expblock}
\end{minipage}

\\[-0.5em]

% ===== Freeport Internship =====
\datecell[-1.6ex]{\small 01.2020 – 06.2020} &
\begin{minipage}[t]{\linewidth}\vspace{0pt}
\begin{expblock}
  \footnotesize
  \setlength{\parindent}{0pt}\setlength{\parskip}{0pt}
    \textbf{\href{https://www.fcx.com}{Freeport-McMoRan Cerro Verde, Arequipa, Peru}}\\[-0.1em]
  \textbf{Maintenance Reliability – Intern}\\[-1em]
  \begin{itemize}[nosep,leftmargin=1em]
    \item Multiphysics CFD simulations (Ansys Fluent) to analyse wear mechanisms and evaluate protective coatings for service-life optimization.
 
  \end{itemize}
\end{expblock}
\end{minipage}

\end{tabular}

\newpage


\resheading{TECHNICAL SKILLS}
\vspace{0.2em}

{\scriptsize
\begin{tabularx}{\textwidth}{@{}X X@{}}

% ====== LEFT COLUMN ======
\begin{minipage}[t]{\linewidth}

\textbf{\textcolor{linksoft}{Programming \& Scientific Computing}}: 
Python (NumPy, Pandas, SciPy, scikit-learn, PyTorch), SQL, C/C++, MATLAB; 
object-oriented design and numerical modeling 
\par\vspace{0.5em}

\textbf{\textcolor{linksoft}{Machine Learning \& AI Systems}}: 
End-to-end ML system design for predictive maintenance and RUL modeling; 
LSTM (Attention), CNN, XGBoost, Random Forest; 
time-series forecasting, anomaly detection; 
Explainable AI (SHAP, PDP) 
\par\vspace{0.5em}

\textbf{\textcolor{linksoft}{Data Engineering \& ETL}}: 
PostgreSQL, Snowflake, InfluxDB (SQL \& NoSQL); 
Dagster, Apache Airflow; 
database schema design, batch and streaming pipelines

\end{minipage}
&
% ====== RIGHT COLUMN ======
\begin{minipage}[t]{\linewidth}

\textbf{\textcolor{linksoft}{Cloud \& Scalable ML Deployment}}: 
Azure ML (pipelines, model registry, deployment, automated retraining), 
Databricks; Kafka, MQTT; scalable cloud-based ML services 
\par\vspace{0.5em}

\textbf{\textcolor{linksoft}{MLOps \& DevOps}}: 
Git, Docker, CI/CD; experiment tracking, 
reproducible environments, containerized inference services 
\par\vspace{0.5em}

\textbf{\textcolor{linksoft}{Industrial \& Physical Systems}}: 
OPC UA integration, Siemens TIA Portal (PLC), 
Bosch Rexroth CtrlX Core; 
FEM (Ansys), CFD; physics-informed modeling 
for cyber-physical systems

\end{minipage}
\\
\end{tabularx}
}


\vspace{0.2em}

\resheading{PUBLICATIONS}\\[\lsepe]
\vspace{0.4em}

% Define lengths for a single column
\newlength{\PubDate}
\newlength{\PubDesc}
\newlength{\PubLink}

\setlength{\PubDate}{1.4cm}
\setlength{\PubLink}{0.4cm}
\setlength{\PubDesc}{\dimexpr(\textwidth - \PubDate - \PubLink - 1.5em)\relax}

\noindent
\begin{tabular}{@{}%
  >{\raggedleft\arraybackslash}m{\PubDate}%
  @{\hspace{0.4em}}%
  >{\RaggedRight\arraybackslash}m{\PubDesc}%
  @{\hspace{0.4em}}%
  >{\centering\arraybackslash}m{\PubLink}@{}}

% Publication 1
\footnotesize 2026 &
\footnotesize \textbf{Explainable Edge-AI Framework for Fault Diagnosis and Remaining Useful Life Prediction in Hydraulic Systems}. 
Published in \emph{Discover Mechanical Engineering} (Springer Nature), 2026.  
DOI: \href{https://doi.org/10.1007/s44245-026-00364-y}{10.1007/s44245-026-00364-y} &
\cert{https://link.springer.com/article/10.1007/s44245-026-00364-y} \\[-0.5em]

\multicolumn{3}{@{}r@{}}{\tiny \textcolor{linksoft}{\tiny[Z] Article}} \\

\end{tabular}
% ========= EDUCATION =========
\resheading{EDUCATION}\\[\lsepe]

\vspace{1em}

% ===== Define lengths =====
\newlength{\colA}
\newlength{\colB}
\newlength{\colC}
\newlength{\colD}

\setlength{\colA}{2.4cm}  % Date
\setlength{\colD}{0.3cm}  % Column for [i], [ii], etc.
\setlength{\colC}{1.5cm}  % Final grade
\setlength{\colB}{\dimexpr\textwidth - \colA - \colC - \colD - 3em\relax} % Main content

% ===== TABLE =====
\noindent
\begin{tabular}{@{}%
  >{\centering\arraybackslash}m{\colA}      % Date (vertically centered, right aligned)
  @{\hspace{1em}}%
  >{\raggedleft\arraybackslash}m{\colD}      % [i], [ii] (vertically centered, right aligned)
  @{\hspace{0.3em}}%
  >{\RaggedRight\arraybackslash}m{\colB}     % Content (vertically centered, left aligned)
  @{\hspace{1em}}%
  >{\raggedleft\arraybackslash}m{\colC}      % Grade (vertically centered, right aligned)
@{}}


%---------- PhD ----------
\footnotesize 06.2026 -- Today & &
\footnotesize \textbf{PhD Candidate – Chair of Sustainable Mobile Drivetrains} &
\footnotesize — \\[-0.15em]
& &
\footnotesize \textit{Topic: Context-Aware Multimodal Learning for Fault Analysis in Software-Intensive Systems.} \\[-0.15em]
& &
\footnotesize \textbf{\href{https://www.tum.de}{Technical University of Munich, Munich, Germany}} & \\[-0.15em]

%---------- MSc 1 ----------
\footnotesize 10.2025 -- 04.2026 &
\raisebox{0.28ex}{\href{https://manuelcaipocc.github.io/cv-caipo-latex/cv_links/uni_ulm_master/Transcript_of_Records_MSc_Computational_Science_Manuel_Caipo.pdf}{\tiny[i]}} &
\footnotesize \textbf{Master's Exchange – Computational Science and Engineering} &
\footnotesize 1,59 \\[-0.15em]
& &

\footnotesize Focus: Numerical Optimization, High-Performance Computing. \\[-0.15em]

& & 
\footnotesize \textbf{ \href{https://www.uni-ulm.de/en/study/study-at-ulm-university/study-programmes/course-information/course/computational-science-and-engineering-master/}{Ulm University, Ulm, Germany }}& \\[-0.15em]

% ---------- MSc 2 ----------
\footnotesize 09.2024 -- 04.2026 &
\raisebox{0.28ex}{\href{https://manuelcaipocc.github.io/cv-caipo-latex/cv_links/hfu_master/Transcript_of_Records_MSc_Advanced_Engineering_Manuel_Caipo.pdf}{\tiny[i]}} &
\footnotesize \textbf{Master of Science – Advanced Precision Engineering} &
\footnotesize 1,59\\[-0.15em] &  &
\footnotesize \textbf{\href{https://www.hs-furtwangen.de/studium/studiengaenge/advanced-precision-engineering-master}{Furtwangen University of Applied Sciences, Villingen-Schwenningen, Germany}}& \\[-0.15em]&
\raisebox{1.55ex} {\href{https://manuelcaipocc.github.io/cv-caipo-latex/cv_links/hfu_master/M_Sc__Thesis_APE_Manuel_Caipo.pdf}{\tiny[ii]}} &
\footnotesize Thesis: Fluid 4.0-compliant framework for digital representation and graph-based neural modeling (published, Springer). & \\[-0.5em]

% ---------- Postgraduate ----------
\footnotesize 06.2021 -- 02.2022 &
\raisebox{0.28ex}{\href{https://siucarne.sunedu.gob.pe/autoridades/informacion-solicitud?Guid=137d2700-4dc5-47dd-8ae4-1192ba5a66b9\&Dato=Z2d6Q0RveUdMTTdJTWtjK2piZ1pBVXUrUjJTV1VWNnVqUFpzR3FST2dlbz0\%3D}{\tiny[i]}} &
\footnotesize \textbf{Postgraduate Diploma – Machine Learning \& Deep Learning} &
\textcolor{mygrey}{\tiny [Top 5\,\%]}   \footnotesize 19.6/20.0\\[-0.6em]& &
\footnotesize \textbf{\href{https://dc.ucsp.edu.pe/postgrado/diplomado-machine-learning/}{Universidad Católica San Pablo, Arequipa, Peru}}& \\[-0.5em]

% ---------- Bachelor ----------
\footnotesize 03.2015 -- 12.2019 &
\raisebox{0.28ex}{\href{https://siucarne.sunedu.gob.pe/autoridades/informacion-solicitud?Guid=f37661dd-9383-4532-8e95-53ce4a8a7fc9\&Dato=aUhGL09TekErTE51eUIvb2VKTERYbE9iVjloaUdFN05xMUVkaStPS3Q0WT0\%3D}{\tiny[i]}}&
\footnotesize \textbf{Bachelor – Mechanical Engineering} &
\textcolor{mygrey}{\tiny [Top 1\,\%]}   \footnotesize 15.2/20.0\\[-0.7em]
& &
\footnotesize \textbf{ \href{https://www.unsa.edu.pe/}{National University of San Agustín, Arequipa, Peru}} & \\[-0.15em]
&\raisebox{0.28ex}{\href{https://manuelcaipocc.github.io/cv-caipo-latex/cv_links/unsa_bachelor/Anerkennung_maschinenbau_manuel_caipo.pdf}{\tiny[ii]}} &
\footnotesize \textit{Official Recognition by the German State} & \\
% & \raisebox{0.28ex}{\href{https://manuelcaipocc.github.io/cv-caipo-latex/cv_links/unsa_bachelor/Beste_note_student.pdf}{\tiny[iii]}} &
% \footnotesize \textit{Award for the highest academic performance.} & \\
&\raisebox{1.55ex}{\href{https://zenodo.org/records/17233982}{\tiny[iv]}} &
\footnotesize Thesis: Design of a 350 TPH copper concentrate conveyor belt with automatic hydraulic tensioning system and loading hopper & \\[-0.5em]

% % ---------- Secondary Education ----------
% \footnotesize 03.2009 -- 12.2013 &
% \raisebox{0.30ex}{\href{https://manuelcaipocc.github.io/cv-caipo-latex/cv_links/school_sekundar/sekundarstufe_HZB.pdf}{\tiny[i]}} &
% \footnotesize \textbf{Secondary Education}&
% \textcolor{mygrey}{\tiny [Top 5\,\%]}    \footnotesize 16.0/20.0 \\[-0.7em]
% &
% &
% \footnotesize \textbf{ \href{https://salesianos.pe/inspectoria/nuestras-obras/arequipa/}{Salesian School Don Bosco }}& \\[-0.15em]

% % ---------- Primary School ----------
% \footnotesize 03.2003 -- 12.2008 & &
% \footnotesize \textbf{Primary School} &
% \footnotesize 18.0/20.0  \\[-0.2em]
% & &
% \footnotesize Charles-Gauss School & \\

\\[-2em]
% ---------- Footnote ----------
\multicolumn{4}{@{}r@{}}{\tiny \vspace{0.5em}
\textcolor{linksoft}{\tiny[i] Diploma}}\\[-0.2em]
% \multicolumn{4}{@{}r@{}}{\tiny 
% \textcolor{linksoft}{\tiny Both Master's programs: Language of instruction German.}}\\




\end{tabular}


\resheading{AWARDS \& HONORS}

% Define lengths for a single column
\newlength{\AwardDate}
\newlength{\AwardDesc}
\newlength{\AwardCert}

\setlength{\AwardDate}{1.4cm}   % date column
\setlength{\AwardCert}{0.4cm}   % certificate column [C]
\setlength{\AwardDesc}{\dimexpr(\textwidth - \AwardDate - \AwardCert - 1.5em)\relax}

\noindent
\begin{tabular}{@{}%
  >{\raggedleft\arraybackslash}m{\AwardDate}%
  @{\hspace{0.4em}}%
  >{\RaggedRight\arraybackslash}m{\AwardDesc}%
  @{\hspace{0.4em}}%
  >{\centering\arraybackslash}m{\AwardCert}@{}}

% Award 1 — Triple Award Freeport
\footnotesize 2022 & 
\footnotesize \textbf{Innova 2022 – Triple Award Winner} (1st Place Digital Transformation, Showroom \& President’s Award) – Freeport-McMoRan.  
\emph{Machine \& Deep Learning Life Data Analysis} with measurable impact (+1.5\% asset availability, ≈>\$10M USD/year). &
\cert{https://manuelcaipocc.github.io/cv-caipo-latex/cv_links/rewards/rewards_innova.pdf} 
\\[0.2em]

% Award 2 — Outstanding Thesis
\footnotesize 2021 & 
\footnotesize \textbf{Outstanding Bachelor's Thesis Award} – Engineering Sciences. &
\cert{https://manuelcaipocc.github.io/cv-caipo-latex/cv_links/rewards/Public_contratulations_thesis.pdf} \\[0.2em]

% Award 3 — Top Student
\footnotesize 2019 & 
\footnotesize \textbf{Top Graduate (Top 1\%)} – Mechanical Engineering, National University of San Agustín, Peru. &
\cert{https://manuelcaipocc.github.io/cv-caipo-latex/cv_links/unsa_bachelor/Beste_note_student.pdf} \\[-0.5em]

\multicolumn{3}{@{}r@{}}{\tiny \textcolor{linksoft}{\textit{Note:} \tiny[Z] Certificate / Letter}} \\

\end{tabular}


% ====== REFERENCES (ATS-optimized) ======
\resheading{REFERENCES}

% % Lengths (if not already defined)
% \newlength{\RefDate}
% \newlength{\RefDesc}
% \newlength{\RefCert}

% \setlength{\RefDate}{1.4cm}   % date column
% \setlength{\RefCert}{0.4cm}   % certificate column [Z]
% \setlength{\RefDesc}{\dimexpr(\textwidth - \RefDate - \RefCert - 1.5em)\relax}

% \noindent
% \begin{tabular}{@{}%
%   >{\raggedleft\arraybackslash}m{\RefDate}%
%   @{\hspace{0.4em}}%
%   >{\RaggedRight\arraybackslash}m{\RefDesc}%
%   @{\hspace{0.4em}}%
%   >{\centering\arraybackslash}m{\RefCert}@{}}

% % Ref 1 — PhD Recommendation
% \footnotesize 2025 & 
% \footnotesize \textbf{Prof. Dr.-Ing. G. Ketterer}, Furtwangen University of Applied Sciences, Dean of Studies M.Sc. \emph{Advanced Precision Engineering}.  
% Recommendation for PhD studies. & 
% \cert{https://manuelcaipocc.github.io/cv-caipo-latex/cv_links/recommendations/Herr_Ketterer_recommendation.pdf} \\[0.2em]

% % Ref 2 — Career / Advanced Studies
% \footnotesize 2023 & 
% \footnotesize \textbf{T. Gaddie}, Freeport-McMoRan, Data Scientist III, Metals Optimization.  
% Reference for career development \& advanced studies. & 
% \cert{https://manuelcaipocc.github.io/cv-caipo-latex/cv_links/recommendations/Empfehlungsschreiben_travis_gaddie.pdf} \\[0.2em]

% % Ref 3 — Master Recommendation
% \footnotesize 2021 & 
% \footnotesize \textbf{Dr. L. Rodríguez B.}, UNSA, Director of Mechanical Engineering \& Bachelor's Thesis Supervisor.  
% Recommendation for Master's studies. & 
% \cert{https://manuelcaipocc.github.io/cv-caipo-latex/cv_links/recommendations/recommendation_unsa.pdf}
% \\[-0.5em]

% \multicolumn{3}{@{}r@{}}{\tiny \textcolor{linksoft}{\textit{Note:} \tiny[Z] Letter of Recommendation / Reference Letter}} \\
% \end{tabular}



% Lengths (if not already defined)
\newlength{\RefDate}
\newlength{\RefDesc}
\newlength{\RefCert}

\setlength{\RefDate}{1.4cm}
\setlength{\RefCert}{0.4cm}
\setlength{\RefDesc}{\dimexpr(\textwidth - \RefDate - \RefCert - 1.5em)\relax}

\noindent
\begin{tabular}{@{}%
  >{\raggedleft\arraybackslash}m{\RefDate}%
  @{\hspace{0.4em}}%
  >{\RaggedRight\arraybackslash}m{\RefDesc}%
  @{\hspace{0.4em}}%
  >{\centering\arraybackslash}m{\RefCert}@{}}

% Ref 1 — Thesis Advisor
\footnotesize 2026 &
\footnotesize \textbf{Dr. Jörn Kretschmer, M.Sc.}, Hochschule Furtwangen University.  
Supervisor of my Master's thesis — Recommendation for PhD studies \href{mailto:Joern.Kretschmer@hs-furtwangen.de}{Joern.Kretschmer@hs-furtwangen.de}. &
\cert{https://manuelcaipocc.github.io/cv-caipo-latex/cv_links/recommendations/Dr_Kretschmer_Letterrecommendation_MCaipo.pdf}
\\[0.2em]

% Ref 2 — PhD Recommendation
\footnotesize 2025 &
\footnotesize \textbf{Prof. Dr.-Ing. G. Ketterer}, Hochschule Furtwangen University, Dean of Studies M.Sc. \emph{Advanced Precision Engineering}.  
Recommendation for PhD studies — \href{mailto:gunter.ketterer@hs-furtwangen.de}{gunter.ketterer@hs-furtwangen.de}. &
\cert{https://manuelcaipocc.github.io/cv-caipo-latex/cv_links/recommendations/Herr_Ketterer_recommendation.pdf}
\\[0.2em]

% Ref 3 — Career / Industry
\footnotesize 2023 &
\footnotesize \textbf{Travis Gaddie}, Freeport-McMoRan, Data Scientist III, Metals Optimization.  
Reference for career development and advanced studies — \href{mailto:tgaddie@fmi.com}{tgaddie@fmi.com}. &
\cert{https://manuelcaipocc.github.io/cv-caipo-latex/cv_links/recommendations/Empfehlungsschreiben_travis_gaddie.pdf}
\\[0.2em]

% Ref 4 — Bachelor Supervisor
\footnotesize 2021 &
\footnotesize \textbf{Dr. L. Rodríguez B.}, National University of San Agustín (UNSA), Director of Mechanical Engineering \& Bachelor's thesis supervisor.  
Recommendation for Master's studies — \href{mailto:lrodriguez4@unsa.edu.pe}{lrodriguez4@unsa.edu.pe}. &
\cert{https://manuelcaipocc.github.io/cv-caipo-latex/cv_links/recommendations/recommendation_unsa.pdf}
\\[-0.5em]

\multicolumn{3}{@{}r@{}}{\tiny \textcolor{linksoft}{\textit{Note:} [Z] Reference letter}} \\
\end{tabular}

% Below AWARDS & HONORS
\resheading{SCHOLARSHIPS \& GRANTS}

% Lengths (if not already defined)
\newlength{\RefDatea}
\newlength{\RefDesca}
\newlength{\RefCerta}

\setlength{\RefDatea}{1.4cm}   % date column
\setlength{\RefCerta}{0.4cm}   % certificate column [C]
\setlength{\RefDesca}{\dimexpr(\textwidth - \RefDate - \RefCert - 1.5em)\relax}

\noindent
\begin{tabular}{@{}%
  >{\raggedleft\arraybackslash}m{\RefDate}%
  @{\hspace{0.4em}}%
  >{\RaggedRight\arraybackslash}m{\RefDesc}%
  @{\hspace{0.4em}}%
  >{\centering\arraybackslash}m{\RefCert}@{}}

% Scholarship 2 — SERESSA
\footnotesize 2019 & 
\footnotesize \textbf{Research Grant – SERESSA 2019}, Spain. 
15th International School on Radiation Effects. & \\[0.2em]

% Scholarship 1 — Bachelor
\footnotesize 2017 & 
\footnotesize \textbf{President of the Republic Scholarship – Peru}.  
Full scholarship awarded for academic excellence (Top 1\%). & 

\end{tabular}

\resheading{ADDITIONAL COURSES}

% ================== GLOBAL LENGTHS ==================
% Separators and spacing control
\newlength{\ColSep}      % space between internal columns
\newlength{\MidSep}      % space between the two halves (between the 2 cells of the outer table)
\newlength{\RowTight}    % vertical compaction per row

% NEW: Proportions for each half
\newlength{\LeftWidth}   % total width of the left table
\newlength{\RightWidth}  % total width of the right table

% Column widths of the LEFT table (date, course, certificate)
\newlength{\Ai}   % date (left)
\newlength{\Bi}   % course (left)
\newlength{\Ci}   % cert   (left)

% Column widths of the RIGHT table (date, course, certificate)
\newlength{\Aii}  % date (right)
\newlength{\Bii}  % course (right)
\newlength{\Cii}  % cert   (right)

% ================== DEFAULT VALUES (ADJUST AS NEEDED) ==================
\setlength{\ColSep}{0.4em}       % internal separation between columns
\setlength{\MidSep}{0.8em}       % separation between the two halves
\setlength{\RowTight}{0.1em}     % vertical compaction per row

% HERE YOU DEFINE THE PROPORTIONS: left 40%, right 60%
\setlength{\LeftWidth}{0.45\textwidth}
\setlength{\RightWidth}{\dimexpr\textwidth - \LeftWidth - \MidSep\relax}

% Left-side widths
\setlength{\Ai}{0.9cm}
\setlength{\Ci}{0.35cm}
\setlength{\Bi}{\dimexpr \LeftWidth - \Ai - \Ci - 2\ColSep \relax}

% Right-side widths
\setlength{\Aii}{\Ai}
\setlength{\Cii}{\Ci}
\setlength{\Bii}{\dimexpr \RightWidth - \Aii - \Cii - 2\ColSep \relax}

% ================== OUTER TABLE (1 ROW x 2 COLUMNS) ==================
\noindent
\begin{tabular}{@{}p{\LeftWidth}@{\hspace{\MidSep}}p{\RightWidth}@{}}
% ================== LEFT CELL ==================
\begin{minipage}[t]{\linewidth}\vspace{0pt}
  \scriptsize
  \begin{tabular}{@{}%
    >{\raggedleft\arraybackslash}m{\Ai}%
    @{\hspace{\ColSep}}%
    >{\RaggedRight\arraybackslash}m{\Bi}%
    @{\hspace{\ColSep}}%
    >{\centering\arraybackslash}m{\Ci}@{}}

  % -------- LEFT-SIDE ROWS --------
  2024 & German Courses B2, C1 (\inst{https://www.vh-ulm.de/}{Ulm vhs} \& \inst{https://www.berlitz.de}{Berlitz}) & \cert{https://manuelcaipocc.github.io/cv-caipo-latex/cv_links/german/testdaf-digital-certificate-scale.pdf} \\[\RowTight]
  2022 & Leadership and Soft Skills (\inst{https://www.ucsp.edu.pe}{UCSP}) & \cert{https://manuelcaipocc.github.io/cv-caipo-latex/cv_links/courses/Leadership_soft_skills.pdf} \\[\RowTight]
  % 2022 & SQL (\inst{https://www.newhorizons.com}{New Horizons}) & \cert{https://manuelcaipocc.github.io/cv-caipo-latex/cv_links/courses/SQL_new_horizons.pdf} \\[\RowTight]
  % 2021 & Python for Data Science (\inst{https://www.socialdataconsulting.com}{SDC}) & \cert{https://manuelcaipocc.github.io/cv-caipo-latex/cv_links/courses/Python_data_science.pdf} \\[\RowTight]
  2021 & Hydraulic Systems (\inst{https://www.boschrexroth.com}{Bosch Rexroth}) & \cert{https://manuelcaipocc.github.io/cv-caipo-latex/cv_links/courses/Rexroth_certificates.pdf} \\[\RowTight]
  2019 & Pipeline Calculation (\inst{https://www.eadic.com}{EADIC}) & \cert{https://manuelcaipocc.github.io/cv-caipo-latex/cv_links/courses/Design_industrial_plants_eadic.pdf} \\[\RowTight]

  \end{tabular}
\end{minipage}
&
% ================== RIGHT CELL ==================
\begin{minipage}[t]{\linewidth}\vspace{0pt}
  \scriptsize
  \begin{tabular}{@{}%
    >{\raggedleft\arraybackslash}m{\Aii}%
    @{\hspace{\ColSep}}%
    >{\RaggedRight\arraybackslash}m{\Bii}%
    @{\hspace{\ColSep}}%
    >{\centering\arraybackslash}m{\Cii}@{}}

  % -------- RIGHT-SIDE ROWS --------
  2020 & Structural \& Multiphysics Analysis with FEM (\inst{https://www.esss.co}{ESSS}) & \cert{https://manuelcaipocc.github.io/cv-caipo-latex/cv_links/courses/ANSYS_Courses.pdf} \\[\RowTight]
  2019 & Industrial Plant Systems (\inst{https://www.tecsup.edu.pe}{TECSUP}) & \cert{https://academico-cloud.tecsup.edu.pe/pcc/\#/home/certificado?c=350995\&n=AQP\%2FPCC\%2F17\%2F7778\&t=E} \\[\RowTight]
  2019 & Autodesk Nastran – FEM \& CFD (\inst{https://www.autodesk.com/de}{Autodesk}) & \cert{https://manuelcaipocc.github.io/cv-caipo-latex/cv_links/courses/Inventor_nastran.pdf} \\[\RowTight]
  2019 & Autodesk Inventor – CAD (\inst{https://www.autodesk.com/de}{Autodesk}) & \cert{https://manuelcaipocc.github.io/cv-caipo-latex/cv_links/courses/Inventor_certificate.pdf} \\[\RowTight]
  2018 & Process Improvement (\inst{https://www.proavance.com}{Proavance}) & \cert{https://manuelcaipocc.github.io/cv-caipo-latex/cv_links/courses/Process_management_improvement.pdf} \\[\RowTight]
      &                                                           & \\[-1em]
  \multicolumn{3}{@{}r@{}}{\tiny \textcolor{linksoft}{[Z] Certificate}}\\
  \end{tabular}
\end{minipage}
\\ % end of the single row of the outer table
\end{tabular}


\end{document}
